\part{複素解析}
	\chapter{複素数}
		\section{複素数の定義}
			任意の2つの実数$a,b\in\realNumbers$に対する2項組$(a,b)$\;(※ここでは(,)という記号を用いているが、どんな記号を用いるかはどうでもいい。別の記号を使っても差し支えない。)を「\textcolor{red}{複素数}」と\underline{名付ける}。
			複素数$(a,b)$の第1要素$a$を「\textcolor{red}{実部}」、第2要素$b$を「\textcolor{red}{虚部}」と\underline{名付ける}。
			複素数全体の集合を$\complexNumbers$と表記することにする。
			\par
			任意の2つの複素数$(a_1,b_1),(a_2,b_2)\in\complexNumbers$に対して演算「和(+)」,「実数倍(スカラー倍)」及び「積($\times$)」を次のように\textcolor{red}{定義}する。
			\begin{align*}
				\textrm{和} &: (a_1,b_1)+(a_2,b_2) = (a_1+a_2,\;b_1+b_2)\\
				\textrm{実数倍}\;(a\in\mathbb{R}) &: a(a_1,b_1) = (aa_1,\;ab_1)\\
				\textrm{積} &: (a_1,b_1)\times(a_2,b_2) = (a_1a_2-b_1b_2,\;a_1b_2+b_2a_1)
			\end{align*}
			\par
			特に、実部が0である複素数$(0,b)$を「\textcolor{red}{純虚数}」と名付ける。
			純虚数のうち虚部が1であるもの$(0,1)$を特に「\textcolor{red}{虚数単位}」と名付け、\textcolor{red}{$i$}と表記することにする。\\
			\par
			$(a,0)$を簡略化して$a$と書くことを許せば、読者が高校以来慣れ親しんでいるであろう、次のような計算が意味を持つ。\\
			例えば
			\[(a_1+ib_1)(a_2+ib_2) = (a_1a_2-b_1b_2)+i(a_1b_2+b_2a_1)\]
			なぜならば、
			\begin{align*}
				(a_1+ib_1)(a_2+ib_2) &= \left[(a_1,0)+i(b_1,0)\right]\times\left[(a_2,0)+i(b_2,0)\right] \\
				&= \left[(a_1,0)+(0,1)\times(b_1,0)\right]\times\left[(a_2,0)+(0,1)\times(b_2,0)\right] \\
				&= \left[(a_1,0)+(0,b_1)\right]\times\left[(a_2,0)+(0,b_2)\right] \\
				&= (a_1,b_1)\times(a_2,b_2) = (a_1a_2-b_1b_2,\;a_1b_2+b_2a_1) \\
				&= (a_1a_2-b_1b_2)+i(a_1b_2+b_2a_1)
			\end{align*}
			また、例えば、
			\[i^2 = -1\]
			なぜならば
			\[i^2 = (0,1)\times(0,1) = (0\times0-1\times1,\;0\times1+1\times0) = (-1,0) = -1\]
			\underline{$i^2 = -1$が虚数単位の定義なのではなく、そうなるように先回りして巧妙に複素数を構築しただけ}であることに注意して欲しい。
			\subsection{注意}
				実数$a$と複素数$(a,0)$は\textcolor{red}{等しくない}。
				イカメシく言うと実数と複素数は「型(type)」が異なる。
				複素数はあくまで上で定義した2実数の\textcolor{red}{組}であって、単一の実数ではない。
				だから、\textcolor{red}{実数と複素数の間}に演算「和」は定義されていない。
				つまり、$a+ib$という書き方は厳密には\textcolor{red}{誤り}である。
				本来$(a,0)+i\times(b,0)$であるものを簡略化して書いているに過ぎない。
			\subsection{コラム \texorpdfstring{$i^2=-1\;?$}{虚数単位の2乗は-1?}}
				「2乗して\textcolor{red}{実数}-1になる」数は無い!! 代わりに、「2乗して\textcolor{red}{複素数}(-1,0)になる」数はある。
				それは先程定義した虚数単位$i=(0,1)$である。
				\par
				いい加減な教科書では、「2乗して-1になる数を虚数単位として定義し、$i$と表記する」と言っているが、この定義は\textcolor{red}{間違っている}。
				複素数の集合の定義が完了していない段階で演算「2乗」と口走っているのがおかしい。
				\par
				\underline{まず数の集合があって、その上に初めて演算を定義できる。}
				先述の「2乗して...」の「2乗」はこの文章中では未定義なのである!!
				百歩譲って実数の集合$\mathbb{R}$上で定義された「2乗」と解釈したとしても、実数の集合に属さない数に対してこの演算を適用できる保証など無い。
				ガソリンの給油口に紅茶を注いで「走れ」と祈るようなものだ。
				\par
				高校で「2乗して-1になる\textcolor{red}{仮想的}な数(\textcolor{red}{imaginary} number)を「虚数単位」$i$と呼ぶ。」
				と教わったはいいものの、「そんな数は自然界にないよなぁ。なんか気持ち悪い...」という\textcolor{red}{違和感}に取り憑かれる原因は、そもそもの定義の間違いにある。
				違和感があって当然である。
				\par
				確かに、人類が複素数を考えた動機は「2乗して-1になる数はないか？」という問題であった。
				新しい理論がそういう動機から始まるのは自然である。
				しかし数学的に「ちゃんとした」理論を作るためには、既存の枠組みのままでは不可能だと分かり次第、思い切ってルールを拡張し、所望の性質が得られるように理屈を先回りして演算を機械的に定義するような break through が必要なこともある。
				その際に、既存の理論に矛盾しないように構築しなければならないのは言うまでもない。
		\section{諸公式}
			\subsection{\texorpdfstring{$x>0,\;z\in\complexNumbers$}{x>0, zは複素数}のとき\texorpdfstring{$|x^z|=x^\Re{z}$}{x\string^zの絶対値はx\string^Re(z)}}
				\begin{proof}
					\quad\par
					$z = a + bi\;(a,b\in\realNumbers)$とする。
					\[ |x^z| = |x^{a+bi}| = |x^a x^{bi}| = |x^a| |x^bi| = |x^\Re{z}| |e^{bi\log{x}}| = x^\Re{z} \]
				\end{proof}
			\subsection{\texorpdfstring{$^\forall n\nmid k,\;\sum_{l=1}^n{\exp\left(i\frac{k}{n}2\pi l\right)} = 0$}{kはnの倍数でないとする。exp(ikl2π/n)をl=1,2...,nについて足すと0である。}}
				\label{fml:花びら定理}
				\begin{proof}
					\begin{align*}
						\sum_{l=1}^n{\exp\left(i\frac{k}{n}2\pi l\right)} &= \frac{\exp\left(i\frac{k}{n}2\pi\right)\left(1-\exp\left(i\frac{k}{n}2\pi\right)^n\right)}{1-\exp\left(i\frac{k}{n}2\pi\right)} = \frac{\exp\left(i\frac{k}{n}2\pi\right)\left(1-e^{ik2\pi}\right)}{1-\exp\left(i\frac{k}{n}2\pi\right)} =0
					\end{align*}
				\end{proof}
		\section{代数方程式}
			\subsection{諸定理}
				\subsubsection{実係数\texorpdfstring{$n$}{n}次方程式が複素数解\texorpdfstring{$s=\alpha+i\beta$}{s=α+iβ}を解に持つならばその複素共役\texorpdfstring{$s^* = \alpha-i\beta$}{}もまた解である}
					\begin{shadebox}
						実係数$n$次方程式が複素数解$s=\alpha+i\beta$を解に持つならばその複素共役$s^* = \alpha-i\beta$もまた解である。
					\end{shadebox}
					\begin{proof}
						\quad\par
						最高次の係数を常に1として一般性を失わないので方程式は
						\[x^n + \sum_{k=n-1}^{0}{c_ix^k} = 0,\quad c_k \in \realNumbers\]
						と書ける。$s=\alpha+i\beta,\;\;\alpha,\beta\in\realNumbers$がこれの解であるから当然
						\[s^n + \sum_{k=n-1}^{0}{c_ks^k} = 0\]
						両辺の複素共役をとると
						\begin{align*}
							\left[s^n + \sum_{k=n-1}^{0}{c_ks^k}\right]^* &= 0^* = 0\\
							\therefore \; {s^*}^n + \sum_{k=n-1}^{0}{c_k{s^*}^k} &= 0
						\end{align*}
						これは$s^*$が解であることに他ならない。
					\end{proof}
			\subsection{例題}
				\subsection{\texorpdfstring{$z^n+z^{n-1}+\dots+z^2+z+1 = 0$}{zのべき乗の和=0}}
					\begin{shadebox}
						次の方程式を解け。
						\[z^n+z^{n-1}+\dots+z^2+z+1 = 0\]
					\end{shadebox}
					\begin{solution}
						\quad\par
						$z=1$が解ではないから両辺に$z-1$を掛けて
						\[z^{n+1}-1=0\]
						の形にする。$z$の絶対値が1だと分かるので$z=e^{\theta i}\; (0\leq\theta \leq 2\pi)$とおいて解き、$z=1$という不要な解を除去する。
						\qed
					\end{solution}
	\chapter{複素引数凸関数}
		\section{複素引数凸関数の例}
			\label{複素引数凸関数の例}
			\subsection{\texorpdfstring{$\|\sum_{i=1}^N X_i\bm{a}_i + \bm{b}\|_2^2\;(\bm{a}_i, \bm{b} \in \complexNumbers^n,\;X_i \in \complexNumbers^{m\times n})$}{行列X1,X2,...,XNの線形結合×ベクトル+定数のノルム}は\texorpdfstring{$X_1,\dotsc,X_N$}{X1,X2,...,XN}に関して凸である。}
				\begin{proof}
					\quad\par
					$T_X \coloneqq (X_1,\dotsc,X_N),\;T_Y \coloneqq (Y_1,\dotsc,Y_N),\;\lambda \in [0,1]$とし、$f(T_X) \coloneqq \|\sum_{i=1}^N X_i\bm{a}_i + \bm{b}\|_2^2$とする。
					\begin{align*}
						f((1-\lambda)T_X + \lambda T_Y) &= \left\|\sum_{i=1}^N ((1-\lambda)X_i + \lambda Y_i)\bm{a}_i + \bm{b}\right\|_2^2 \\
						&= \left\|(1-\lambda)\left(\sum_{i=1}^N X_i\bm{a}_i + \bm{b}\right) + \lambda\left(\sum_{i=1}^N Y_i\bm{a}_i + \bm{b}\right)\right\|_2^2 \\
						&= (1-\lambda)^2\left\|\sum_{i=1}^N X_i\bm{a}_i + \bm{b}\right\|_2^2 + \lambda^2\left\|\sum_{i=1}^N Y_i\bm{a}_i + \bm{b}\right\|_2^2 \\
						&\phantom{=} + 2\lambda(1-\lambda)\Re{\left(\sum_{i=1}^N X_i\bm{a}_i + \bm{b}\right)^\mathrm{H}\left(\sum_{i=1}^N Y_i\bm{a}_i + \bm{b}\right)}
					\end{align*}
					よって
					\begin{align*}
						&\phantom{=} (1-\lambda)f(T_X) + \lambda f(T_Y) - f((1-\lambda)T_X + \lambda T_Y) \\
						&= \lambda(1-\lambda)\left[\left\|\sum_{i=1}^N X_i\bm{a}_i + \bm{b}\right\|_2^2 + \left\|\sum_{i=1}^N Y_i\bm{a}_i + \bm{b}\right\|_2^2 - 2\Re{\left(\sum_{i=1}^N X_i\bm{a}_i + \bm{b}\right)^\mathrm{H}\left(\sum_{i=1}^N Y_i\bm{a}_i + \bm{b}\right)}\right] \\
						&= \lambda(1-\lambda)\left\|\left(\sum_{i=1}^N X_i\bm{a}_i + \bm{b}\right) - \left(\sum_{i=1}^N Y_i\bm{a}_i + \bm{b}\right)\right\|_2^2 \geq 0
					\end{align*}
				\end{proof}
			\subsection{\texorpdfstring{$\norm{A\bm{x}+\bm{b}}_2^2$}{||Ax+b||\string^2}が狭義凸\texorpdfstring{$\iff A^*A\succ O$}{⇔ A*Aが正定}}
				\begin{shadebox}
					$A \in \complexNumbers^{m\times n}, \bm{x} \in \complexNumbers^n, \bm{b} \in \complexNumbers^m$とする。
					$f(\bm{x}) \coloneqq \norm{A\bm{x}+\bm{b}}_2^2$が$\bm{x}$に関して狭義凸となる必要十分条件は$A^*A\succ O$である。
				\end{shadebox}
				$\lambda \in (0,1),\; \bm{x},\bm{y} \in \complexNumbers^m, \bm{x}\neq\bm{y}$とする。
				\begin{proof}
					まず次式が成り立つ。
					\begin{align*}
						f((1-\lambda)\bm{x} + \lambda\bm{y}) &\coloneqq \norm{(1-\lambda)(A\bm{x}+\bm{b}) + \lambda(A\bm{y}+\bm{b})}_2^2 \\
						&= (1-\lambda)^2 f(\bm{x}) + \lambda^2 f(\bm{y}) + 2\lambda(1-\lambda)\Re{(A\bm{x}+\bm{b})^*(A\bm{y}+\bm{b})}
					\end{align*}
					これより
					\begin{align*}
						&\phantom{=} (1-\lambda)f(\bm{x}) + \lambda f(\bm{y}) - f((1-\lambda)\bm{x} + \lambda\bm{y}) \\
						&= \lambda(1-\lambda)\left[f(\bm{x}) + f(\bm{y}) - 2\Re{(A\bm{x}+\bm{b})^*(A\bm{y}+\bm{b})}\right] \\
						&= \lambda(1-\lambda)\norm{(A\bm{x}+\bm{b}) - (A\bm{y}+\bm{b})}_2^2 = \lambda(1-\lambda)\norm{A(\bm{x}-\bm{y})}_2^2 \\
						&= \lambda(1-\lambda)(\bm{x}-\bm{y})^*A^*A(\bm{x}-\bm{y})
					\end{align*}
					$\lambda \in (0,1),\;\bm{x}\neq\bm{y}$に対して上式が正となる必要十分条件は$A^*A\succ O$である。
				\end{proof}
	\chapter{特殊関数}
		\section{Gamma関数}
			\subsection{\texorpdfstring{$\lim_{n\to\infty} \integrate{0}{n}{t^{z-1}(1+t/n)^n}{}{t} = \Gamma(z)\;(z\in\complexNumbers, \Re{z} > 0)$}{Gamma関数の無限乗積表示}}
				Wikipediaで紹介されている無限乗積表示の導出過程でこの主張が使われているが、その根拠が薄かったので自力で証明してみた。
				\begin{proof}
					\quad\par
					$\varepsilon>0$を任意にとる。
					$G_n(z) \coloneqq \integrate{0}{n}{t^{z-1}(1+t/n)^n}{}{t}$とする。
					\par
					Gamma関数$\integrate{0}{\infty}{t^{z-1}e^{-t}}{}{t}$は収束するから、十分大きい$T_1(\varepsilon)>0$を取れば
					\[ \left|\integrate{T_1(\varepsilon)}{\infty}{t^{z-1}e^{-t}}{}{t}\right| \leq \integrate{T_1(\varepsilon)}{\infty}{\left|t^{z-1}e^{-t}\right|}{}{t} < \varepsilon/3 \]
					\par
					\ref{(1-t/n)^nは単調増加}より$(1-t/n)^n$は$n$について単調増加であり、かつ$\lim_{n\to\infty}(1-t/n)^n = e^{-t}$なので、$n \geq T_1(\varepsilon)$ならば
					\begin{align*}
						&\left|\integrate{T_1(\varepsilon)}{n}{t^{z-1}(1-t/n)^n}{}{t}\right| \leq \integrate{T_1(\varepsilon)}{n}{\left|t^{z-1}(1-t/n)^n\right|}{}{t} \\
						&= \integrate{T_1(\varepsilon)}{n}{t^{\Re{z}-1}(1-t/n)^n}{}{t} \leq \integrate{T_1(\varepsilon)}{\infty}{\left|t^{z-1}e^{-t}\right|}{}{t} < \varepsilon/3
					\end{align*}
					また、十分大きい自然数$N_1(\varepsilon,T_1(\varepsilon))$をとると、$n \geq N_1(\varepsilon,T_1(\varepsilon))$ならば次式が成り立つ。
					\[ \left|\integrate{0}{T_1(\varepsilon)}{t^{z-1}\{e^{-t} - (1-t/n)^n\}}{}{t}\right| < \varepsilon/3 \]
					以上より、$n \geq \max\{T_1(\varepsilon), N_1(\varepsilon,T_1(\varepsilon))\}$ならば次式が成り立つ。
					\begin{align*}
						\left|\Gamma(z) - G_n(z)\right| &= \left|\integrate{T_1(\varepsilon)}{\infty}{t^{z-1}e^{-t}}{}{t} + \integrate{0}{T_1(\varepsilon)}{t^{z-1}\{e^{-t} - (1-t/n)^n\}}{}{t} - \integrate{T_1(\varepsilon)}{n}{t^{z-1}(1-t/n)^n}{}{t}\right| \\
						&\leq \left|\integrate{T_1(\varepsilon)}{\infty}{t^{z-1}e^{-t}}{}{t}\right| + \left|\integrate{0}{T_1(\varepsilon)}{t^{z-1}\{e^{-t} - (1-t/n)^n\}}{}{t}\right| + \left|\integrate{T_1(\varepsilon)}{n}{t^{z-1}(1-t/n)^n}{}{t}\right| < \varepsilon
					\end{align*}
				\end{proof}
	\chapter{複素積分}
		\section{例題}
			\subsection{\texorpdfstring{$\integrate{0}{2\pi}{\Log(1-a\exp(i\theta))}{}{\theta}=0\quad (C\text{: 原点中心の単位円},\;a\in\realNumbers)$}{Log(1-a exp(iΘ))の単位円周上での積分は0}}
				\begin{solution}
					\quad\par
					$a=0$のときは明らかに成り立つ。$a\neq0$のときは$z=e^{i\theta}$とすると
					\[\integrate{0}{2\pi}{\Log(1-ae^{i\theta})}{}{\theta} = \integrate{C}{}{\frac{\Log(1-az)}{iz}}{}{z}\quad (C:\text{原点中心の左回り単位円})\]
					被積分関数の$C$内部の特異点は$z=0$である。原点の近傍(具体的には$|z|<1/|a|$の範囲)においては
					\[\frac{\Log(1-az)}{iz} = \frac{-1}{iz}\sum_{n=1}^\infty\frac{(az)^n}{n}=\frac{-1}{i}\sum_{n=1}^\infty\frac{a^nz^{n-1}}{n} = \frac{-1}{i}\sum_{n=0}^\infty\frac{a^{n+1}}{n+1}z^n\]
					であり、特異点$z=0$は除去可能であることから従う。
				\end{solution}
	\chapter{Fourier級数}
		\section{連続関数のFourier級数の高周波成分は0に収束する}
			\label{連続関数のFourier級数の高周波成分は0に収束する}
			\begin{shadebox}
				$f$は区間$[-\pi,\pi]$で連続とする。$f$のFourier係数を$c_k$とすると$\lim_{k\to\infty}c_k = 0$となる。
			\end{shadebox}
			\begin{proof}
				\quad\par
				まず
				\[ \lim_{k\to\infty}\integrate{-\pi}{\pi}{f(x)\sin kx}{}{x} = 0 \]
				を示す。ここで
				\[ \integrate{-\pi}{\pi}{f(x)\sin kx}{}{x} = \integrate{0}{\pi}{(f(x)-f(-x))\sin kx}{}{x} \]
				であり、$f(x)-f(-x)$は$[0,\pi]$で連続なので結局のところ区間$[0,\pi]$の任意の連続関数$f$に対して
				\[ \lim_{k\to\infty}\integrate{0}{\pi}{f(x)\sin kx}{}{x} = 0 \]
				が成り立つことを示せば良い。
				$f$は閉区間で連続だから有界であり、特に一様連続なので、任意の$\varepsilon>0$に対してある$\delta>0$が存在して
				$\forall (x,y\in[0,\pi],|x-y|\leq\delta),|f(x)-f(y)|<\varepsilon$が成り立つ。
				$k \geq 2,\;\pi/k < \delta$となるように$k$を十分大きくとる。
				\begin{align*}
					\integrate{0}{\pi}{f(x)\sin kx}{}{x} &= \sum_{l=0}^{\floor{k/2}-1}\left(\integrate{2l\pi/k}{(2l+1)\pi/k}{f(x) \sin kx}{}{x} + \integrate{(2l+1)\pi/k}{(2l+2)\pi/k}{f(x) \sin kx}{}{x}\right) \\
					&\quad + \integrate{2\floor{k/2}\pi/k}{\pi}{f(x)\sin kx}{}{x}
				\end{align*}
				最後の項は$k\to\infty$のとき$0$に収束する。
				最初の2項は$y=kx$と変数変換して
				\begin{align*}
					&\quad \sum_{l=0}^{\floor{k/2}-1}\frac{1}{k}\left(\integrate{2l\pi}{(2l+1)\pi}{f(y/k) \sin y}{}{y} + \integrate{(2l+1)\pi}{(2l+2)\pi}{f(y/k) \sin y}{}{y}\right) \\
					&= \frac{1}{k} \sum_{l=0}^{\floor{k/2}-1} \integrate{2l\pi}{(2l+1)\pi}{(f(y/k) - f(y/k+\pi/k)) \sin y}{}{y}
				\end{align*}
				となる。これの絶対値は
				\begin{align*}
					&\leq \frac{1}{k} \sum_{l=0}^{\floor{k/2}-1} \integrate{2l\pi}{(2l+1)\pi}{\abs{f(y/k) - f(y/k+\pi/k) \sin y}}{}{y} \\
					&< \frac{1}{k} \sum_{l=0}^{\floor{k/2}-1} \integrate{2l\pi}{(2l+1)\pi}{\varepsilon\abs{\sin y}}{}{y} = \frac{2\varepsilon}{k} \sum_{l=0}^{\floor{k/2}-1}1 = \frac{2\varepsilon}{k}\floor{k/2} \leq \frac{2\varepsilon}{k}\frac{k}{2} = \varepsilon
				\end{align*}
				同様にして
				\[ \lim_{k\to\infty}\integrate{-\pi}{\pi}{f(x)\cos kx}{}{x} = 0 \]
				も示せる。
			\end{proof}
		\section{\texorpdfstring{$\integrate{-\pi}{\pi}{\frac{\sin nx}{2\tan\frac{x}{2}}}{}{x} = \pi\;(n\in\naturalNumbers)$}{sin(nx)/(2 tan(x/2))の[-π,π]での積分はπ}}
			\label{integ{-pi}{pi}{(sin nx)/2tan(x/2)}{x} = pi}
			\begin{proof}
				\quad\par
				求めたい積分の値を$I_n$とする。
				分子の$\sin nx$について$nx = (n-1/2)x + x/2$として加法定理により展開すると
				\[ I_n = \frac{1}{2}\integrate{-\pi}{\pi}{\frac{\cos^2(x/2)}{\sin(x/2)}\sin(n-1/2)x}{}{x} + \frac{1}{2}\integrate{-\pi}{\pi}{\cos(x/2)\cos(n-1/2)x}{}{x} \]
				第一項の分子で$\cos^2(x/2) = 1-\sin^2(x/2)$とすると
				\begin{align*}
					I_n &= \frac{1}{2}\integrate{-\pi}{\pi}{\frac{\sin(n-1/2)x}{\sin(x/2)}}{}{x} - \frac{1}{2}\integrate{-\pi}{\pi}{\sin(x/2)\sin(n-1/2)x}{}{x} + \frac{1}{2}\integrate{-\pi}{\pi}{\cos(x/2)\cos(n-1/2)x}{}{x} \\
					&= \frac{1}{2}\integrate{-\pi}{\pi}{\frac{\sin(n-1/2)x}{\sin(x/2)}}{}{x} + \frac{1}{2}\integrate{-\pi}{\pi}{\cos nx}{}{x} \quad (\text{↑の第2,3項を加法定理でまとめた}) \\
					&= \frac{1}{2}\integrate{-\pi}{\pi}{\frac{\sin(n-1/2)x}{\sin(x/2)}}{}{x} \quad (\because\;\text{第2項は}0)
				\end{align*}
				さらに$(n-1/2)x = (n-1)x + x/2$として加法定理を使うと
				\[
					I_n = \frac{1}{2}\integrate{-\pi}{\pi}{\frac{\sin(n-1)x\;\cos(x/2)}{\sin(x/2)}}{}{x} + \frac{1}{2}\integrate{-\pi}{\pi}{\cos(n-1)x}{}{x} = I_{n-1} +
					\begin{cases}
						\pi & (n=1) \\
						0 & (\text{otherwise})
					\end{cases}
				\]
				これと$I_0 = 0$より定理が成り立つ。
			\end{proof}
		\section{Fourier級数の応用}
			\subsection{\texorpdfstring{$\integrate{0}{\infty}{\frac{\sin x}{x}}{}{x} = \pi/2$}{sin(x)/xの0から∞までの積分はπ/2}}
			\begin{proof}
				\quad\par
				複素積分を使った導出がよく知られているが、実解析の範囲で導出してみる。
				この積分が収束することは、$(\sin x)/x$が$[0,\pi/2]$で連続であることから$[0,\pi/2]$で積分可能であること、および$\sin x = (-\cos x)'$を用いて部分積分することで$[\pi/2,\infty)$で積分が絶対収束することから示せる。
				よって、
				\begin{align*}
					\integrate{0}{\infty}{\frac{\sin x}{x}}{}{x} &= \lim_{n\to\infty}\integrate{0}{n\pi}{\frac{\sin x}{x}}{}{x} = \lim_{n\to\infty}\integrate{0}{\pi}{\frac{\sin nt}{t}}{}{t} \quad (t=x/n\text{とおいた}) \\
					&= \lim_{n\to\infty}\frac{1}{2}\integrate{-\pi}{\pi}{\frac{\sin nt}{t}}{}{t} = \lim_{n\to\infty}\frac{1}{2}\integrate{-\pi}{\pi}{\left(\frac{1}{t} - \frac{1}{2\tan(t/2)} + \frac{1}{2\tan(t/2)}\right)\sin nt}{}{t}
				\end{align*}
				$\frac{1}{t} - \frac{1}{2\tan(t/2)}$は$[-\pi,\pi]$で連続だから\ref{連続関数のFourier級数の高周波成分は0に収束する}より
				\[ \lim_{n\to\infty}\frac{1}{2}\integrate{-\pi}{\pi}{\left(\frac{1}{t} - \frac{1}{2\tan(t/2)}\right)\sin nt}{}{t} = 0 \]
				であること、および\ref{integ{-pi}{pi}{(sin nx)/2tan(x/2)}{x} = pi}より証明が済む。
			\end{proof}
	\chapter{Fourier変換}
		\section{余弦,正弦変換}
		{%スコープ。\ge上書き回避
			\renewcommand{\ge}{g_\textrm{e}}
			\newcommand{\go}{g_\textrm{o}}
			\newcommand{\Ge}{G_\textrm{e}}
			\newcommand{\Go}{G_\textrm{o}}
			$f(t)$の余弦変換$C(\omega)$,正弦変換$S(\omega)$は次のように定義される。
			\begin{align*}
				C(\omega) &\coloneqq \sqrt{\frac{2}{\pi}}\integrate{0}{\infty}{f(x)\cos{\omega t}}{}{t}\\
				S(\omega) &\coloneqq \sqrt{\frac{2}{\pi}}\integrate{0}{\infty}{f(x)\sin{\omega t}}{}{t}
			\end{align*}
			$C,S$はそれぞれ\textcolor{red}{$\omega$の偶関数,奇関数}になっていることに注意。そして$t\geq0$において
			\begin{equation}
				f(t) \sim \sqrt{\frac{2}{\pi}}\integrate{0}{\infty}{C(\omega)\cos{\omega t}}{}{\omega} = \sqrt{\frac{2}{\pi}}\integrate{0}{\infty}{S(\omega)\sin{\omega t}}{}{\omega} \nonumber
			\end{equation}
			が成り立つ。注意すべきは、この逆変換は$t<0$に関しては何も教えてくれないということだ。変換時に$t<0$における情報を捨ててしまっているのだから当然である。
			以下、上の関係を示す。
			\newline\par
			\begin{proof}
				\quad\par
				$f(t)$の$t\geq0$の部分を取り出して$t<0$の部分に拡張することを考える。関数$\ge(t),\go(t)$を次のように定義する。
				\[\ge(t) \coloneqq f(|t|)\;\;(-\infty<t<\infty)\]
				\begin{align}
					\go(t) \coloneqq
					\begin{cases}
						f(t) & (t>=0) \\
						-f(-t) & (t<0)
					\end{cases}
					\nonumber
				\end{align}
				$\ge,\go$はそれぞれ$f$を偶関数,奇関数として拡張したものである。両者のフーリエ変換を$\Ge(\omega),\Go(\omega)$とすると
				\begin{align*}
					\Ge(\omega) &= \frac{1}{\sqrt{2\pi}}\integrate{-\infty}{\infty}{\ge(t)e^{-i\omega t}}{}{t} = \frac{1}{\sqrt{2\pi}}\integrate{-\infty}{\infty}{\ge(t)(\cos{\omega t}-i\sin{\omega t})}{}{t}\\
					&= \frac{1}{\sqrt{2\pi}}\integrate{-\infty}{\infty}{\ge(t)\cos{\omega t}}{}{t}\;\;\;(\because \ge(t)\sin{\omega t}\text{は奇関数なので対称積分の結果は0})\\
					&= \sqrt{\frac{2}{\pi}}\integrate{0}{\infty}{\ge(t)\cos{\omega t}}{}{t} = \sqrt{\frac{2}{\pi}}\integrate{0}{\infty}{f(t)\cos{\omega t}}{}{t} = C(\omega)
				\end{align*}
				となり余弦変換と一致する。同様にして
				\[\Go(\omega) = -iS(\omega)\]
				となることも確かめられる。
				\newline\par
				次に逆フーリエ変換で$\ge,\go$を復元してみる。
				\begin{align*}
					\ge(t) &\sim \frac{1}{\sqrt{2\pi}}\integrate{-\infty}{\infty}{\Ge(\omega)e^{i\omega t}}{}{\omega} = \frac{1}{\sqrt{2\pi}}\integrate{-\infty}{\infty}{C(\omega)e^{i\omega t}}{}{\omega} = \frac{1}{\sqrt{2\pi}}\integrate{-\infty}{\infty}{C(\omega)(\cos{\omega t}+i\sin{\omega t})}{}{\omega}\\
					&= \frac{1}{\sqrt{2\pi}}\integrate{-\infty}{\infty}{C(\omega)\cos{\omega t}}{}{\omega}\;\;(\because C(\omega)\sin{\omega t}\text{ は奇関数なので対称積分の結果は0})\\
					&= \sqrt{\frac{2}{\pi}}\integrate{0}{\infty}{C(\omega)\cos{\omega t}}{}{\omega}
				\end{align*}
				\begin{align*}
					\go(t) &\sim \frac{1}{\sqrt{2\pi}}\integrate{-\infty}{\infty}{\Go(\omega)e^{i\omega t}}{}{\omega} = \frac{1}{\sqrt{2\pi}}\integrate{-\infty}{\infty}{-iS(\omega)e^{i\omega t}}{}{\omega} = \frac{1}{\sqrt{2\pi}}\integrate{-\infty}{\infty}{-iS(\omega)(\cos{\omega t}+i\sin{\omega t})}{}{\omega}\\
					&= \frac{1}{\sqrt{2\pi}}\integrate{-\infty}{\infty}{S(\omega)\sin{\omega t}}{}{\omega}\;\;(\because S(\omega)\cos{\omega t}\text{ は奇関数なので対称積分の結果は0})\\
					&= \sqrt{\frac{2}{\pi}}\integrate{0}{\infty}{S(\omega)\sin{\omega t}}{}{\omega}
				\end{align*}
				$t\geq0$では$f,\ge,\go$は等しいから
				\[f(t) \sim \sqrt{\frac{2}{\pi}}\integrate{0}{\infty}{C(\omega)\cos{\omega t}}{}{\omega} = \sqrt{\frac{2}{\pi}}\integrate{0}{\infty}{S(\omega)\sin{\omega t}}{}{\omega}\]
			\end{proof}
			繰り返しになるが上式は$t\geq0$でしか成り立たない。$t<0$のことはわからないのである。クレイジーな変換だと思うかもしれないが、ちゃんと利点もある。実験データや音声,映像データは普通、記録開始を時刻0として扱う。この場合は$t<0$に関しては知る必要がない(どうでもいい)のだ。そういう状況下では複素数が登場する本家のフーリエ変換よりもsin,cosどっちか好きな方だけ使えば済むような正弦,余弦変換が役に立つことがある。
		}
	\chapter{Bessel関数}
		\section{定義}
			$x\in\realNumbers,\;z\in\complexNumbers$に対して$e^{iz\sin x}$のFourier級数展開により「$n$次の第一種Bessel関数$J_n(z)$」を次のように定義する。
			\begin{equation}
				\label{eqn:implicit definition of 1st kind n-th Bessel function}
				e^{iz\sin x} = \sum_{n = -\infty}^\infty J_n(z)e^{inx}
			\end{equation}
			すなわち
			\begin{equation}
				\label{eqn:explicit definition of 1st kind n-th Bessel function}
				J_n(z) = \frac{1}{2\pi}\integrate{-\pi}{\pi}{e^{iz\sin x}e^{-inx}}{}{x}
			\end{equation}
			\subsection{簡単化}
				式\eqref{eqn:explicit definition of 1st kind n-th Bessel function}は次のようにして簡単化できる。
				\[ J_n(z) = \frac{1}{2\pi}\integrate{-\pi}{\pi}{\exp i\left(z\sin x -nx\right)}{}{x} = \frac{1}{2\pi}\integrate{-\pi}{\pi}{\cos\left(z\sin x -nx\right)}{}{x} + \frac{1}{2\pi}\integrate{-\pi}{\pi}{i\sin\left(z\sin x -nx\right)}{}{x} \]
				第1項は偶関数の対称積分であり、第2項は奇関数の対称積分なので0だから結局次のようになる。
				\begin{equation}
					\label{eqn:simplified definition of 1st kind n-th Bessel function}
					J_n(z) = \frac{1}{2\pi}\integrate{-\pi}{\pi}{\cos\left(z\sin x -nx\right)}{}{x} = \frac{1}{\pi}\integrate{0}{\pi}{\cos\left(z\sin x -nx\right)}{}{x}
				\end{equation}
		\section{性質}
			\subsection{\texorpdfstring{$J_n(-z) = J_{-n}(z)$}{y軸に関する鏡像反転}}
				式\eqref{eqn:simplified definition of 1st kind n-th Bessel function}より容易に示せる。
			\subsection{\texorpdfstring{$J_{2m}(-z) = J_{2m}(z),\;J_{2m+1}(-z) = -J_{2m+1}(z)$}{次数と偶関数,奇関数}}
				$m\in\integers$とするとき、次式が成り立つ。
				\[ J_{2m}(-z) = J_{2m}(z),\;J_{2m+1}(-z) = -J_{2m+1}(z) \]
				\begin{proof}
					\quad\par
					$n=2m$とする。
					\begin{align*}
						J_{2m}(-z) &= \frac{1}{2\pi}\integrate{-\pi}{\pi}{\cos(-z\sin x - 2mx)}{}{x} \\
						&= \frac{1}{2\pi}\integrate{-2\pi}{0}{\cos(z\sin y - 2m(y+\pi))}{}{y} \quad (\text{変数変換 }x=y+\pi) \\
						&= \frac{1}{2\pi}\integrate{-\pi}{\pi}{\cos(z\sin y - 2m(y+\pi))}{}{y} = \frac{1}{2\pi}\integrate{-\pi}{\pi}{\cos(z\sin y - 2my)}{}{y} = J_{2m}(z)
					\end{align*}
					同様にして$J_{2m+1}(-z) = -J_{2m+1}(z)$も示せる。
				\end{proof}
		\section{Maclaurin展開}
			式\eqref{eqn:implicit definition of 1st kind n-th Bessel function}の左辺を指数関数に関してMaclaurin展開すると$x$を変数としたFourier級数展開が得られることを利用して$J_n(z)$のMaclaurin展開を求める。
			\[ e^{iz\sin x} = \exp\left( z\frac{e^{ix} - e^{-ix}}{2} \right) = \sum_{k=0}^\infty \frac{1}{k!} \left(\frac{z}{2}\right)^k (e^{ix}-e^{-ix})^k \] %= \sum_{k=0}^\infty \frac{1}{k!} \left(\frac{z}{2}\right)^k \sum_{l=0}^k \binom{k}{l} (-1)^{k-l} e^{i(2l-k)x}
			$n\geq 0$のとき、$(e^{ix}-e^{-ix})^k$の展開で$e^{inx}$が生じるのは$k=n+2l\;(l=0,1,2,\dotsc)$のときで、その係数は$(-1)^l\binom{k}{l} = (-1)^l\binom{n+2l}{l}$である。
			$n<0$のとき、$(e^{ix}-e^{-ix})^k$の展開で$e^{inx}$が生じるのは$k=-n+2l\;(l=0,1,2,\dotsc)$のときで、その係数は$(-1)^{k-l}\binom{k}{l} = (-1)^{-n+l}\binom{-n+2l}{l}$である。
			式\eqref{eqn:implicit definition of 1st kind n-th Bessel function}の両辺のFourier展開基底の係数同士が一致せねばならぬから
			\[
				J_n(z) =
				\begin{dcases}
					\sum_{l=0}^\infty \frac{(-1)^l}{(n+2l)!}\left(\frac{z}{2}\right)^{n+2l}\binom{n+2l}{l} = \sum_{l=0}^\infty \frac{(-1)^l}{l!(n+l)!}\left(\frac{z}{2}\right)^{n+2l} & (n\geq 0) \\
					\sum_{l=0}^\infty \frac{(-1)^{-n+l}}{(-n+2l)!}\left(\frac{z}{2}\right)^{-n+2l}\binom{-n+2l}{l} = \sum_{l=0}^\infty \frac{(-1)^{-n+l}}{l!(-n+l)!}\left(\frac{z}{2}\right)^{-n+2l} & (n < 0)
				\end{dcases}
			\]
